\chapter{Método}
\label{cha:metodo}

\section{Delineamento da pesquisa}

Este estudo adota uma revisão bibliográfica estruturada, de caráter qualitativo e comparativo, orientada pelas recomendações do PRISMA 2020 e pela extensão PRISMA-S para relato das estratégias de busca. O objetivo é identificar como a literatura científica recente define o metaverso, quais características são recorrentes, quais tecnologias aparecem como habilitadoras e quais aplicações são descritas. A revisão é utilizada como base para confrontar a definição adotada no projeto e construir uma definição operacional em camadas.

O PRISMA 2020 orienta o relato transparente das fontes de informação, critérios de elegibilidade, processo de seleção e síntese. O PRISMA-S acrescenta requisitos específicos para que a busca seja reproduzível, incluindo o registro da fonte, da data e da estratégia completa utilizada em cada fonte \cite{page2021,rethlefsen2021}.

\section{Questões de pesquisa}

A busca foi estruturada para responder às seguintes questões:

\begin{itemize}
    \item \textbf{QP1:} Quais características aparecem com maior recorrência nas definições científicas de metaverso?
    \item \textbf{QP2:} Quais características da definição adotada neste projeto são sustentadas pela literatura?
    \item \textbf{QP3:} Quais características podem ser unificadas sem ampliar excessivamente o conceito?
    \item \textbf{QP4:} Quais elementos devem ser tratados como tecnologias habilitadoras, e não como requisitos conceituais?
    \item \textbf{QP5:} Quais lacunas permanecem na definição adotada e quais implicações elas trazem para o projeto?
\end{itemize}

\section{Período e fontes de informação}

O protocolo foi definido para contemplar publicações em inglês entre janeiro de 2022 e setembro de 2026, período escolhido para concentrar a análise na literatura produzida após a consolidação do debate acadêmico contemporâneo sobre o metaverso.

As fontes planejadas são:

\begin{itemize}
    \item Scopus;
    \item Web of Science;
    \item IEEE Xplore;
    \item ACM Digital Library;
    \item Google Scholar, como fonte complementar de recuperação.
\end{itemize}

A escolha é coerente com revisões anteriores do tema. Tukur et al. utilizaram IEEE Xplore, Scopus, ACM e Google Scholar e reportaram 447 registros após restrições de busca; Ritterbusch e Teichmann utilizaram Web of Science, JSTOR, Wiley e AIS, com 381 registros iniciais; e Boztas et al. realizaram uma revisão de definições e arquiteturas a partir do Google Scholar, iniciando com 237 registros \cite{tukur2024,ritterbusch2023,boztas2025}. Esses números pertencem às respectivas revisões e não são tratados como resultados desta pesquisa.

\section{Estratégia de busca}

Foi definida uma string conceitual principal:

\begin{verbatim}
("metaverse" OR "metaverses")
AND
(definition OR conceptualization OR characteristics OR architecture)
AND
(technology OR technologies OR application OR applications
OR interoperability OR persistence)
\end{verbatim}

Também foi definida uma string específica para localização de revisões:

\begin{verbatim}
("metaverse" OR "metaverses")
AND
("systematic literature review" OR "systematic review"
OR "scoping review" OR "literature review")
AND
(definition OR characteristics OR technologies OR applications)
\end{verbatim}

A sintaxe deve ser adaptada aos operadores e campos específicos de cada base, mantendo a mesma lógica conceitual. De acordo com o PRISMA-S, a versão efetivamente executada em cada base deve ser preservada exatamente como utilizada, juntamente com a data e os filtros aplicados \cite{rethlefsen2021}.

\section{Critérios de inclusão}

Foram definidos previamente os seguintes critérios:

\begin{itemize}
    \item publicação entre janeiro de 2022 e setembro de 2026;
    \item publicação em inglês;
    \item artigo científico ou trabalho de conferência com contribuição acadêmica identificável;
    \item relação direta com o conceito de metaverso;
    \item discussão de pelo menos uma das dimensões de interesse: definição, características, arquitetura, tecnologias, aplicações, persistência ou interoperabilidade;
    \item texto completo disponível para avaliação;
    \item contribuição suficiente para responder pelo menos uma questão de pesquisa.
\end{itemize}

\section{Critérios de exclusão}

Foram definidos os seguintes critérios de exclusão:

\begin{itemize}
    \item notícias, páginas comerciais, textos de opinião e materiais sem caráter científico;
    \item trabalhos que apenas mencionem o termo metaverso sem analisá-lo;
    \item registros duplicados;
    \item trabalhos sem texto completo disponível;
    \item trabalhos sem relação suficiente com as questões de pesquisa;
    \item estudos cujo foco seja exclusivamente uma tecnologia isolada, sem relação analítica com o metaverso;
    \item documentos fora do período ou idioma definidos.
\end{itemize}

\section{Execução e auditoria da busca}

A busca bibliográfica foi executada em 29 de setembro de 2026 por meio das fontes científicas e páginas públicas acessíveis durante a pesquisa. Foram localizados e verificados registros diretamente relacionados à definição, às características, às arquiteturas, às tecnologias e às aplicações do metaverso.

Entretanto, os contadores ao vivo de resultados das interfaces fechadas de Scopus, Web of Science, IEEE Xplore, ACM Digital Library e Google Scholar não são expostos de maneira confiável no ambiente utilizado nesta pesquisa. Por essa razão, este trabalho não atribui números inventados a essas bases. Os quantitativos finais do fluxograma PRISMA somente poderão ser declarados como resultados desta revisão após a execução das strings nas interfaces das bases e o registro dos respectivos contadores.

Essa decisão metodológica é deliberada: um número de resultados não observado diretamente não deve ser apresentado como dado empírico da revisão.

Como controle externo de plausibilidade, foram comparados os procedimentos e quantitativos publicados por revisões anteriores. Tukur et al. reportaram 212 registros provenientes do IEEE Xplore, 96 do Scopus, 39 da ACM e 100 do Google Scholar, totalizando 447 registros após restrições específicas; 141 duplicatas foram removidas e 122 estudos chegaram à avaliação de texto completo \cite{tukur2024}. Ritterbusch e Teichmann reportaram 381 registros iniciais em Web of Science, JSTOR, Wiley e AIS, chegando a 28 definições ou descrições após as etapas de seleção e buscas de citações \cite{ritterbusch2023}. Boztas et al. identificaram inicialmente 237 registros no Google Scholar, removeram 6 duplicatas e 7 registros por outros motivos, avaliaram 192 textos completos e incluíram 103 estudos \cite{boztas2025}.

\section{Resultados bibliográficos verificados}

A busca realizada permitiu identificar um núcleo atualizado de trabalhos particularmente relevantes para este estudo:

\begin{table}[H]
\centering
\caption{Trabalhos de revisão e definição identificados na busca bibliográfica}
\label{tab:trabalhos_verificados}
\begin{tabular}{p{0.28\textwidth}p{0.23\textwidth}p{0.39\textwidth}}
\toprule
\textbf{Trabalho} & \textbf{Ano} & \textbf{Contribuição para este estudo}\\
\midrule
Ritterbusch e Teichmann & 2023 & Definições e características recorrentes do metaverso.\\
Tukur et al. & 2024 & Técnicas, tecnologias e aplicações do metaverso.\\
Zhou, Chen e Jin & 2024 & Definição, características técnicas e comportamento do usuário em perspectiva sociotécnica.\\
Bénaben, Congès e Fertier & 2025 & Proposta de definição unificada e framework para propriedades do metaverso.\\
Enamorado-Díaz et al. & 2025 & Aplicações, tecnologias e desafios de engenharia de software.\\
Boztas, Ghadafi e Ibrahim & 2025 & Definições e arquiteturas, com análise temática de 33 definições.\\
\bottomrule
\end{tabular}
\end{table}

Os trabalhos identificados reforçam uma convergência importante para o objetivo deste artigo: não há uma definição única e universalmente aceita, mas existem conjuntos recorrentes de características, tecnologias habilitadoras e dimensões sociais, econômicas e arquitetônicas \cite{ritterbusch2023,zhou2024,boztas2025}.

\section{Fluxo PRISMA}

O fluxograma PRISMA definitivo não é preenchido com estimativas. Nesta versão, a tabela abaixo funciona como registro metodológico do estado da busca e separa claramente dados efetivamente observados de dados que ainda dependem da execução manual nas interfaces das bases.

\begin{table}[H]
\centering
\caption{Estado do registro PRISMA}
\label{tab:prisma}
\begin{tabular}{p{0.62\textwidth}p{0.23\textwidth}}
\toprule
\textbf{Etapa} & \textbf{Situação}\\
\midrule
Bases e fontes definidas & Registrado\\
Strings conceituais definidas & Registrado\\
Critérios de inclusão definidos & Registrado\\
Critérios de exclusão definidos & Registrado\\
Busca bibliográfica exploratória executada & Registrado\\
Registros recuperados em cada base fechada & A registrar na interface da base\\
Duplicatas & A calcular após exportação\\
Triagem de títulos e resumos & A executar sobre o conjunto exportado\\
Elegibilidade por texto completo & A executar sobre o conjunto exportado\\
Estudos incluídos & A calcular após a elegibilidade\\
\bottomrule
\end{tabular}
\end{table}

Portanto, esta versão não apresenta o quadro como um PRISMA numérico completo. O preenchimento definitivo depende dos contadores e dos registros exportados das bases. Essa distinção preserva a rastreabilidade e evita que números de revisões anteriores sejam confundidos com resultados desta pesquisa.

\section{Procedimento de análise}

Após a seleção definitiva, cada estudo será classificado segundo cinco dimensões: (i) definição do metaverso; (ii) características estruturais; (iii) tecnologias habilitadoras; (iv) aplicações e setores; e (v) desafios ou limitações.

A síntese será realizada por comparação temática. Características conceituais, como ambiente compartilhado, interação, identidade digital, persistência e imersão, serão separadas das tecnologias que podem viabilizá-las, como realidade estendida, inteligência artificial, computação distribuída e tecnologias descentralizadas. Essa separação evita transformar uma tecnologia específica em requisito universal do conceito.

A etapa final consistirá na comparação entre a definição adotada no projeto e os padrões encontrados na literatura. A definição resultante será apresentada como uma definição operacional para o projeto, e não como uma afirmação de que existe uma única definição científica definitiva.
